\subsection{Spherical Decomposition of Isotropic Fields}\label{subsec:ps_spherical}

This section bridges Cartesian Fourier and spherical coordinates, explaining how an isotropic field decomposes in both bases.

\subsubsection{Completeness of spherical harmonics}

Any function on the unit sphere can be expanded in spherical harmonics:

\begin{equation}
f(\theta, \phi) = \sum_{\ell=0}^\infty \sum_{m=-\ell}^\ell f_{\ell m} Y_{\ell m}(\theta, \phi)
\end{equation}

For isotropic fields, the power depends only on $\ell$:

\begin{equation}
\langle f_{\ell m} f^*_{\ell' m'} \rangle = C_\ell \delta_{\ell\ell'} \delta_{mm'}
\end{equation}

This $\ell$-only dependence is the key to reducing 3D power to angular power spectra.

\subsubsection{Spherical Fourier-Bessel (SFB) basis}

For 3D fields in spherical coordinates $(r, \theta, \phi)$, the natural orthonormal basis combines radial and angular:

\begin{equation}
\text{Basis function: } j_\ell(kr) Y_{\ell m}(\theta, \phi)
\end{equation}

where $j_\ell$ are spherical Bessel functions and $k$ is the radial wavenumber.

Any field expands as:

\begin{equation}
f(\mathbf{r}) = \sqrt{\frac{2}{\pi}} \int_0^\infty dk \sum_{\ell m} f_{\ell m}(k) k j_\ell(kr) Y_{\ell m}(\theta, \phi)
\label{eq:sfb_expansion_ref}
\end{equation}

See \cref{ch:sfb} for complete treatment of Bessel functions and SFB orthonormality.
